% Estilo LaTeX do artigo de lesões periapicais.
% Adaptado do modelo do artigo ROCKFACE-LCCMAT (segmentação de poros).
% Compilar com LuaLaTeX a partir da pasta layout_latex.
\usepackage[
  paperwidth=210mm,
  paperheight=279mm,
  left=14mm,
  right=12mm,
  top=17mm,
  bottom=15mm,
  headheight=9mm,
  headsep=4mm,
  footskip=8mm
]{geometry}

\usepackage{fontspec}
\defaultfontfeatures{Ligatures=TeX,Scale=MatchLowercase}
\setmainfont[
  Path=../fonts/,
  UprightFont=texgyretermes-regular.otf,
  BoldFont=texgyretermes-bold.otf,
  ItalicFont=texgyretermes-italic.otf,
  BoldItalicFont=texgyretermes-bolditalic.otf
]{TeX Gyre Termes}
\setsansfont[
  Path=../fonts/,
  UprightFont=texgyreheros-regular.otf,
  BoldFont=texgyreheros-bold.otf,
  ItalicFont=texgyreheros-italic.otf,
  BoldItalicFont=texgyreheros-bolditalic.otf
]{TeX Gyre Heros}

\usepackage[portuguese]{babel}
\usepackage{microtype}
\usepackage{graphicx}
\graphicspath{{../figuras/}{../../Resultados/figuras/R/}}
\usepackage{booktabs}
\usepackage{array}
\usepackage{tabularx}
\usepackage{ragged2e}
\usepackage{enumitem}
\usepackage{titlesec}
\usepackage{caption}
\usepackage{xcolor}
\usepackage{fancyhdr}
\usepackage{hyperref}
\usepackage{dblfloatfix}
\usepackage{balance}

\definecolor{rockblue}{HTML}{0065A3}
\definecolor{rockgreen}{HTML}{CAD85E}
\definecolor{rockorange}{HTML}{C45A00}
\definecolor{scopegray}{HTML}{F3F3F3}

\hypersetup{
  colorlinks=true,
  linkcolor=rockblue,
  urlcolor=rockblue,
  citecolor=rockblue,
  pdftitle={Detecção de lesões periapicais em radiografias panorâmicas},
  pdfsubject={Random Forest por pixel versus U-Net; figuras geradas em R e Python},
  pdfauthor={William Jota},
  pdfcreator={LuaLaTeX - ML Odonto}
}
\urlstyle{same}

\setlength{\columnsep}{4mm}
\setlength{\parindent}{4.5mm}
\setlength{\parskip}{0pt}
\emergencystretch=1.2em
\clubpenalty=10000
\widowpenalty=10000
\displaywidowpenalty=10000

\AtBeginDocument{\fontsize{8.97pt}{10.71pt}\selectfont}

\titleformat{\section}
  {\sffamily\bfseries\fontsize{10.21pt}{12pt}\selectfont}
  {\thesection.}{0.35em}{}
\titlespacing*{\section}{0pt}{4mm}{0.8mm}
\titleformat{\subsection}
  {\sffamily\bfseries\fontsize{8.72pt}{10.4pt}\selectfont}
  {\thesubsection.}{0.35em}{}
\titlespacing*{\subsection}{0pt}{1.8mm}{0.3mm}
\titleformat{name=\section,numberless}
  {\sffamily\bfseries\fontsize{10.21pt}{12pt}\selectfont}
  {}{0pt}{}

\captionsetup{
  font={small},
  labelfont=bf,
  labelsep=space,
  justification=raggedright,
  singlelinecheck=false,
  skip=1.7mm
}
% babel (português) redefine os nomes no início do documento; reaplicar depois dele
\addto\captionsportuguese{\renewcommand{\figurename}{Fig.}\renewcommand{\tablename}{Tabela}}

% Mais figuras de largura total por página, para evitar páginas só com uma figura
\renewcommand{\dbltopfraction}{0.92}
\renewcommand{\dblfloatpagefraction}{0.72}
\renewcommand{\textfraction}{0.06}
\setcounter{dbltopnumber}{3}

\pagestyle{fancy}
\fancyhf{}
\fancyhead[L]{\sffamily\fontsize{7.97pt}{9pt}\selectfont\color{rockblue}Draft manuscript · not peer reviewed}
\fancyhead[R]{\sffamily\bfseries\fontsize{9.96pt}{11pt}\selectfont Article}
\fancyfoot[L]{\sffamily\fontsize{7.97pt}{9pt}\selectfont\color{rockblue}ML Odonto | draft of 27 October 2026}
\fancyfoot[R]{\sffamily\bfseries\fontsize{7.97pt}{9pt}\selectfont\thepage}
\renewcommand{\headrulewidth}{0.6pt}
\renewcommand{\footrulewidth}{0.6pt}

\fancypagestyle{firstpage}{
  \fancyhf{}
  \fancyhead[L]{\sffamily\bfseries\fontsize{16.94pt}{19pt}\selectfont manuscript draft}
  \fancyhead[R]{\sffamily\bfseries\fontsize{14.94pt}{17pt}\selectfont Article}
  \fancyfoot[L]{\sffamily\fontsize{7.97pt}{9pt}\selectfont\color{rockblue}ML Odonto | draft of 27 October 2026}
  \fancyfoot[R]{\sffamily\bfseries\fontsize{7.97pt}{9pt}\selectfont\thepage}
  \renewcommand{\headrulewidth}{3.2pt}
  \renewcommand{\headrule}{\hbox to\headwidth{\color{rockgreen}\leaders\hrule height \headrulewidth\hfill}}
  \renewcommand{\footrulewidth}{0.6pt}
}

% #1 data, #2 título, #3 autores, #4 resumo, #5 palavras-chave, #6 afiliação
\newcommand{\articlefront}[6]{%
  \begin{minipage}{\textwidth}
    \vspace{4mm}
    \begin{flushright}
      {\sffamily\fontsize{7.97pt}{9pt}\selectfont\color{rockblue}#1}
    \end{flushright}
    \vspace{-1mm}\hrule\vspace{6mm}
    {\sffamily\bfseries\fontsize{25.9pt}{28pt}\selectfont #2\par}
    \vspace{9mm}
    {\sffamily\bfseries\fontsize{9.22pt}{13pt}\selectfont #3\par}
    \vspace{3mm}\hrule\vspace{2mm}
    {\sffamily\fontsize{9.96pt}{14pt}\selectfont\RaggedRight #4\par}
    \vspace{2.5mm}
    {\sffamily\fontsize{7.47pt}{9.4pt}\selectfont #5\par}
    \vspace{3mm}\hrule\vspace{2mm}
    {\sffamily\fontsize{7.47pt}{9.4pt}\selectfont #6\par}
    \vspace{3mm}
  \end{minipage}
}

\newcommand{\widefigure}[4][!t]{%
  \begin{figure*}[#1]
    \centering
    \includegraphics[width=\textwidth,height=#2,keepaspectratio]{#3}
    \caption{#4}
  \end{figure*}
}

\newcommand{\scopenote}[1]{%
  \par\noindent
    \colorbox{scopegray}{%
      \parbox{\dimexpr\columnwidth-2\fboxsep\relax}{%
        \sffamily\color{black!72}\setlength{\parindent}{0pt}%
        \textcolor{rockorange}{\rule{2pt}{2.6em}}\hspace{2mm}#1}}
  \par\vspace{2mm}
}

% Marcador visível de informação que os autores precisam fornecer
\newcommand{\pendente}[1]{\textcolor{rockorange}{\textbf{[#1]}}}

\newcolumntype{Y}{>{\RaggedLeft\arraybackslash}X}
\newcolumntype{L}{>{\RaggedRight\arraybackslash}X}
\newcommand{\tablefont}{\sffamily\fontsize{7.2pt}{8.6pt}\selectfont}
\setlist{nosep,leftmargin=4mm}
